\section{Prior explicit bounds and the even-power comparison}
\label{sec:literature}

\subsection{An older theorem already gives a number for every power}

The distinction between an explicitly tabulated bound and a bound
that follows by substitution into a general theorem matters here.
Nesterenko--Waldschmidt~\cite{NesterenkoWaldschmidt}, Theorem~2(1),
prove that if \(\xi\) is real algebraic of degree \(d\) and
coefficient length at most \(\mathcal L\ge3\), then
\begin{equation}\label{eq:NW}
|\pi-\xi|\ge
\exp\{-1200000d(\log\mathcal L+d\log d)(1+\log d)\}.
\end{equation}
The coefficient length of an integral polynomial is the sum of the
absolute values of its coefficients; for an algebraic number it is
the length of its primitive integral minimal polynomial. The
following specialization is derived here from their theorem.

\begin{proposition}[A classical all-power consequence]\label{prop:NWpowers}
For every integer \(k\ge1\), put
\[
B(k)=1200000k(1+\log k).
\]
Then \(\mu(\pi^k)\le B(k)\). In particular, explicit exponent
upper bounds for all fixed powers were already available from
classical quantitative transcendence theory.
\end{proposition}
\begin{proof}
It suffices first to consider \(q\ge1\) and
\(|\pi^k-p/q|<1\). This implies \(p/q>0\). Set
\(\theta=(p/q)^{1/k}\), and let \(f\) be its primitive integral
minimal polynomial, of degree \(d\le k\) and positive leading
coefficient \(a\). Gauss's lemma gives
\(f\mid qX^k-p\) in \(\ZZ[X]\), after removing the content of
the latter polynomial if necessary. In particular \(a\le q\).
Every conjugate of \(\theta\) is a root of this binomial, hence
has modulus \(\theta\). Also \(\theta<\pi+1\). The elementary
bound for the coefficients of a product gives
\[
\mathcal L(f)\le a(1+\theta)^d<6^kq.
\]
This step uses the minimal polynomial and remains valid when the
binomial is reducible.

Use \(\mathcal L=6^kq\) in \eqref{eq:NW}. Monotonicity in
\(1\le d\le k\) yields
\[
|\pi-\theta|\ge
\exp\{-B(k)(\log q+k\log6+k\log k)\}
=(6k)^{-kB(k)}q^{-B(k)}.
\]
Finally,
\[
|\pi^k-p/q|=|\pi-\theta|
 \sum_{a=0}^{k-1}\pi^{k-1-a}\theta^a
\ge\pi^{k-1}|\pi-\theta|.
\]
For the fractions excluded at the start the error is at least one.
Thus the smaller prefactor \((6k)^{-kB(k)}\) gives a lower bound
for all integer \(p\) and \(q\ge1\), proving the exponent claim.
\end{proof}

This consequence is not a claim that the original paper prints a
table of \(\pi^k\) exponents. It shows why absence of such a table
would not establish priority for an explicit number. It also gives
an independent comparison whose growth is only \(O(k\log k)\).

\begin{center}
\begin{tabular}{rrr}
\toprule
\(k\)&Classical consequence \(B(k)\), approximately
 &Bound displayed in these notes\\
\midrule
3&\(7\,555\,004.24\)&\(1107.93\)\\
5&\(15\,656\,627.47\)&\(42\,461\)\\
7&\(24\,745\,645.25\)&\(1\,272\,103\)\\
9&\(34\,530\,025.44\)&\(35\,398\,937\)\\
\bottomrule
\end{tabular}
\end{center}

The first three entries show a substantial improvement over this
particular classical consequence. The fourth does not: the old
general theorem already beats the displayed twelve-layer \(k=9\)
number. None of these comparisons is an exhaustive literature
search for the best bound at the indicated power. They delimit the
contribution that can safely be claimed.

\subsection{Even zeta values give an identification, not an approximation theorem}

Euler's formula, with \(B_{2r}\) the Bernoulli numbers, reads
\[
\zeta(2r)=\frac{(-1)^{r+1}B_{2r}(2\pi)^{2r}}{2(2r)!}
\in\QQ^\times\pi^{2r}.
\]
Therefore \(\mu(\zeta(2r))=\mu(\pi^{2r})\). For example,
\[
\zeta(2)=\pi^2/6,\qquad \zeta(4)=\pi^4/90.
\]
This equality transfers an existing irrationality-exponent estimate
between two names for the same rationally scaled constant. It does
not turn evaluation of the zeta value into a proof of a sharp
irrationality measure. Nor does an upper bound for \(\mu(\pi^2)\)
automatically give the corresponding upper bound for \(\mu(\pi^4)\).
Approximating \((\pi^2)^2\) and approximating \(\pi^2\) by rational
numbers are different approximation problems.

Marcovecchio--Zudilin~\cite{MarcovecchioZudilin} establish the
unconditional specialized estimate
\begin{equation}\label{eq:zeta4benchmark}
\mu(\pi^4)=\mu(\zeta(4))\le12.51085940\ldots.
\end{equation}
Their 2019 preprint, published in 2020, uses a different
hypergeometric representation to control the arithmetic of
approximations to \(\zeta(4)\). The bound does not assume the older
denominator conjecture discussed in their introduction. It is the
specialized benchmark verified for this revision; the present notes
do not claim a complete current classification of all later work.

\subsection{What the same L99 shape would do for the fourth power}

For the accepted cubic rates, direct substitution gives
\[
D-3G=544.69491-3(213.81256)=-96.74277.
\]
Thus the fixed accepted cubic instance is inadmissible for \(k=4\).
This is a failed hypothesis, rather than a failure of the general
transfer theorem. Increasing the common multiplicity scale changes
both the degree and the decay-to-height ratio.

For orientation, keep the base L99 shape
\((16,1,\ldots,1)\), and use the root-free lower rate of
Section~\ref{subsec:l99uniform} with scale \(s=19\). Then
\[
\begin{split}
m&=816,\qquad \underline D_s\approx3630.47940,\\
G_s&\approx1019.35966,\qquad
\underline D_s-3G_s\approx572.40043.
\end{split}
\]
The associated expression is approximately \(6628.69640\).
The accompanying numerical check certifies that it is below
\(6628.70\), using the already proved rational divisor lower rate.
This is a manuscript-level specialization and an Arb comparison,
not an accepted fixed Lean theorem for \(\pi^4\). Its size is
already sufficient to explain the comparison: the uniform method
is far from \eqref{eq:zeta4benchmark}.

Consequently there are two different useful projects. Completing
general-power applications tests and documents a reusable proof
framework. Improving the specialized fourth-power bound requires
competing below \(12.51085940\ldots\), most naturally by sharpening
the arithmetic of direct \(\zeta(4)\) approximations. Extending the
range of a theorem and improving the best numerical constant are
distinct achievements.
